\section{Introduction}

\subsection{The failure mode}

A research process can make progress while its own framing becomes the dominant source of error. More search does not repair a missing representation. A larger tree does not by itself reveal that the wrong parent discipline has been treated as irrelevant. More evidence does not justify rewriting a question when the actual defect is merely an execution failure. The problem studied here is therefore not generic iteration, autonomous science, or long-horizon reasoning. It is the narrower governance problem of preserving an inspectable epistemic state while deciding when evidence is strong enough to authorize a change in the scientific formulation itself.

For a scoped problem, the architecture studied here represents the evolving research state as
\[
X_t=(K_t,W_t,M_t),
\]
where $K_t$ is the provenance-preserving state of knowledge about the target object, $W_t$ is the current model of which domains, representations, mechanisms, and search routes may be relevant, and $M_t$ is the governed research method. The three coordinates are separated because they require different evidence and different authority to change. This engineering separation is in the spirit of Parnas's information-hiding principle: change-prone design decisions should be isolated behind explicit interfaces rather than entangled implicitly \cite{parnas1972}. The revision policy for $K_t$ likewise borrows the minimal-change distinction between expansion, contraction, and revision from the AGM belief-revision tradition \cite{agm1985}; no claim is made here that the operational update rules satisfy the full AGM postulates, which would have to be established separately.

A new source may update $K_t$ without changing $W_t$. A previously unrecognized parent discipline or representation may require a change to $W_t$ while leaving the method intact. A persistent, well-attributed failure can become evidence for a challenger to $M_t$, but method revision is governed separately and is not an operation this paper authorizes.

\subsection{Protected scientific contribution}

Paper I isolates one dominant scientific contribution with four coupled requirements. First, $K/W/M$ are explicit co-evolving state rather than hidden conversational context. Second, discoveries and failures target typed formulation coordinates, so diagnosed responsibility constrains which high-level mutation is admissible. Third, high-level change is gated by lower-level exclusion, same-task counterfactual necessity, and protected-invariant preservation rather than generic ``reflect and retry.'' Fourth, an accepted mutation stales exactly the dependent closure identified by its bound impact rather than allowing prior success certificates to survive silently.
\subsection{Contributions}

Four properties of the mechanism are put forward and tested.

The proposed object is the coupling. Recursive audit, active diagnosis, counterfactual repair, dependency rollback, certificate enforcement, evidence-gated revision, and generic mutation authorization remain donor mechanisms. The governed policy composes them into a protected scientific-escalation contract. Responsibility determines the candidate mutation, a counterfactual tests whether high-level change repairs the same task, protected preservation prevents collateral regression, and dependency impact determines what must reopen.

The archived hidden-shift study tests whether that contract can make reformulation successful and selective. A later faithful-comparator challenge shows that ordered search alone lets an active value-of-information parent match the governed policy on the primary joint criterion. The P1-X successor tests the same control law on heterogeneous revision-responsibility contracts after a donor-complete comparison product receives the underlying mechanisms. An information-equivalent product ties exactly.

The resulting claim is bounded. The paper defines and executes an explicit contract for high-level scientific reformulation. The registered studies establish its finite behaviour, failure modes, and intervention-cost profile. They do not establish comparative mechanism necessity, because a competent ordered-search comparator matches the governed policy on the primary joint criterion.

\subsection{Five-paper boundary}

Paper I owns the reconstruction problem. Open-world discovery and stopping are evaluated in Paper II; cross-domain absorption and the global knowledge portrait in Paper III; authority promotion and protected verification in Paper IV; and method evolution in Paper V. Paper I uses those components as explicit interfaces so that the reformulation result remains attributable to one well-defined scientific object rather than to an umbrella architecture.
What is defended empirically is narrower than any of the four taken alone: it is the composition of protected necessity conditions under which a high-level formulation or search-universe change is admissible at all. Section~\ref{sec:relatedwork} places each ingredient with its nearest neighbour in the literature, and Section~\ref{sec:limitations} states what the composition is and is not evidence for.
